\documentclass[11pt]{article}

\usepackage[T1]{fontenc}
\usepackage{amsmath,amssymb,amsthm}
\usepackage[margin=1in]{geometry}
\usepackage{booktabs}
\usepackage[hidelinks]{hyperref}
\hypersetup{
  pdftitle={Nonexistence of an Order-17 C4-Free Graph with 123 Common-Neighbor Pairs},
  pdfauthor={Xiaojian Zeng}
}

\newtheorem{theorem}{Theorem}
\newtheorem{lemma}[theorem]{Lemma}
\newtheorem{proposition}[theorem]{Proposition}
\newcommand{\codeg}{\operatorname{codeg}}

\title{Nonexistence of an Order-17 $C_4$-Free Graph\\
with 123 Common-Neighbor Pairs}
\author{Xiaojian Zeng\\
PBC School of Finance, Tsinghua University, 100190\\
\texttt{xiaojianzeng@foxmail.com}}
\date{}

\begin{document}
\maketitle

\begin{abstract}
We prove that there is no finite simple graph $G$ on $17$ vertices for
which every two distinct vertices have at most one common neighbor and
$\sum_{v\in V(G)}\binom{d(v)}2=123$.  The codegree condition is equivalent
to $C_4$-freeness.  Previous work excluded the planar case and reduced any
remaining witness to maximum degree between $5$ and $9$.  A complete exact
enumeration then leaves eight degree multisets.  Six are excluded by a
local deficiency inequality, the seventh by parity in the leave of the
common-neighbor relation, and the eighth by incompatible lower and upper
bounds on triangle incidences.  Computation is used only for the finite
degree-multiset reduction; all eight surviving cases are eliminated
symbolically.
\end{abstract}

\noindent\textbf{Keywords.}
Common neighbors; $C_4$-free graph; degree sequence; exact enumeration;
triangle incidence.

\medskip
\noindent\textbf{MSC 2020.} 05C35, 05C07, 05C30.

\section{The problem and the result}

For distinct vertices $x,y$ of a simple graph $G$, write
\[
  \codeg_G(x,y)=|N_G(x)\cap N_G(y)|.
\]
Ascolese, Negrini, Pagani, and Pellegrini introduced the corresponding
codegree-realization problem and gave a complete answer within the class of
planar $C_4$-free graphs~\cite{AscoleseEtAl2026}.  In the version consulted,
their Theorem~5.9 excludes the planar graph at the parameter below,
Corollary~5.8 reduces a remaining witness to the nonplanar range
$5\leq\Delta\leq9$, and the sentence following that corollary conjectures
that this remaining case does not occur.  We prove precisely that
nonexistence assertion.

\begin{theorem}\label{thm:main}
There is no finite simple undirected graph $G$ such that
\[
 |V(G)|=17,\qquad
 \codeg_G(x,y)\leq1\quad(x\ne y),\qquad
 \sum_{v\in V(G)}\binom{d(v)}2=123.
\]
\end{theorem}

The finite part of the proof follows the certificate-aware practice of
making the search domain, integer predicates, and independent checks
explicit.  For comparison, this practice is also used in recent
finite-library graph investigations based on exact integer predicates and
independent verifiers~\cite{ZengEtAl2026Certificates}.  No mathematical
result from that work is used here.

\section{Reduction to eight degree multisets}

Suppose for contradiction that $G$ satisfies the hypotheses of
Theorem~\ref{thm:main}.  Two vertices have two distinct common neighbors if
and only if the four vertices contain a (not necessarily induced) $4$-cycle.
Thus $G$ is $C_4$-free.  Moreover,
\[
 \sum_{v}\binom{d(v)}2
   =\sum_{\{x,y\}\in\binom{V(G)}2}\codeg_G(x,y).
\]
The codegree bound therefore says that exactly $123$ of the
$\binom{17}{2}=136$ vertex pairs have one common neighbor and the other
$13$ have none.

The graph is connected.  Indeed, every pair of vertices in different
components has no common neighbor, while a nontrivial partition of a
$17$-element set has at least $1\cdot16=16$ cross-part pairs.

Put $m=|E(G)|$.  The target equation and the handshaking lemma give
\[
 \sum_v d(v)^2=246+2m.
\]
Cauchy--Schwarz yields
\[
 4m^2\leq17(246+2m).
\]
This inequality already fails at $m=37$, and the difference between its
left- and right-hand sides is strictly increasing for $m\geq37$; hence
$m\leq36$.  The cited planar exclusion and degree reduction now imply that
$G$ is nonplanar and
\begin{equation}\label{eq:delta-range}
 5\leq\Delta(G)\leq9.
\end{equation}

We shall also use the following necessary inequality, which is the
specialization of Lemma~5.3 of~\cite{AscoleseEtAl2026}.  Its short proof is
included to make the reduction transparent.

\begin{lemma}\label{lem:source-ineq}
For every vertex $x$ of a $C_4$-free graph on $17$ vertices,
\begin{equation}\label{eq:source-ineq}
 \binom{17-d(x)}2\geq
 \sum_{y\ne x}\binom{d(y)-1}{2}.
\end{equation}
\end{lemma}

\begin{proof}
Let $A=V(G)\setminus N(x)$, so $|A|=17-d(x)$.  For each $y\ne x$, at most
one neighbor of $y$ lies in $N(x)$; consequently
$|N(y)\cap A|\geq d(y)-1$.  If one pair $\{a,b\}\subseteq A$ lay in both
$\binom{N(y)\cap A}{2}$ and $\binom{N(z)\cap A}{2}$ for distinct $y,z$,
then $y$ and $z$ would have the two common neighbors $a,b$.  Hence these
families of pairs are disjoint as $y$ varies.  Counting them inside
$\binom A2$ proves~\eqref{eq:source-ineq}.
\end{proof}

\begin{lemma}[Exact degree reduction]\label{lem:eight}
The degree multiset of $G$ must be one of the eight rows in
Table~\ref{tab:degrees}.  In the table, $a^b$ means that $a$ occurs with
multiplicity $b$, $h$ is the number of degree-$5$ vertices, and
\[
 S=\sum_{d(v)<5}d(v)(5-d(v)).
\]
\end{lemma}

\begin{table}[ht]
\centering
\caption{The complete list surviving the degree-sequence filters.}
\label{tab:degrees}
\begin{tabular}{@{}cclrr@{}}
\toprule
case & $m$ & degree multiset & $h$ & $S$\\
\midrule
1 & 34 & $5^{12},2^3,1^2$       & 12 & 26\\
2 & 34 & $5^{11},4^2,2,1^3$     & 11 & 26\\
3 & 35 & $5^{11},4,3,2^4$       & 11 & 34\\
4 & 35 & $5^{11},3^4,2,1$       & 11 & 34\\
5 & 35 & $5^{10},4^3,3,2^2,1$   & 10 & 34\\
6 & 35 & $5^9,4^5,3,1^2$        &  9 & 34\\
7 & 36 & $5^9,4^3,3^5$          &  9 & 42\\
8 & 36 & $5^8,4^6,3^2,2$        &  8 & 42\\
\bottomrule
\end{tabular}
\end{table}

\begin{proof}
Connectedness,~\eqref{eq:delta-range}, and $m\leq36$ show that it is enough
to inspect nonincreasing $17$-tuples $(d_1,\ldots,d_{17})$ in
$\{1,\ldots,9\}^{17}$ satisfying
\begin{equation}\label{eq:filters}
 5\leq d_1\leq9,\quad
 \sum_i\binom{d_i}{2}=123,\quad
 \sum_i d_i\ \text{is even and at most }72,
\end{equation}
together with graphicality and~\eqref{eq:source-ineq} for every entry.
Graphicality was tested both with the Erd\H{o}s--Gallai criterion
\cite{ErdosGallai1960} and with the Havel--Hakimi reduction
\cite{Havel1955,Hakimi1962}.

For completeness, an unpruned implementation traverses all
\[
 \binom{9+17-1}{17}=\binom{25}{17}=1{,}081{,}575
\]
multisets using a combinations-with-replacement iterator.  It retains a
tuple exactly when the literal predicates in~\eqref{eq:filters}, the
Erd\H{o}s--Gallai inequalities
\[
 \sum_{i=1}^k d_i\leq k(k-1)+\sum_{i=k+1}^{17}\min(d_i,k)
 \quad(1\leq k\leq17),
\]
and~\eqref{eq:source-ineq} all hold.  The successive counts are $1{,}368$
after the arithmetic predicates, $1{,}356$ after graphicality, and $8$
after~\eqref{eq:source-ineq}, giving Table~\ref{tab:degrees}.

Two separately written enumerators reproduce the list: one uses a pruned
sorted-tuple recursion and a direct Erd\H{o}s--Gallai implementation; the
other enumerates degree multiplicities and uses a direct Havel--Hakimi
implementation.  The unpruned cross-check uses NetworkX~3.4.2
\cite{HagbergEtAl2008} and requires its two graphicality methods to agree.
All predicates and counters use integer arithmetic; no random choices or
numerical tolerances occur.  Thus every positive degree multiset allowed by
the preceding mathematical reductions is covered.
\end{proof}

\section{Local deficiency and triangle bounds}

All rows of Table~\ref{tab:degrees} have maximum degree $5$.  Put
\[
 H=\{v:d(v)=5\},\qquad L=V(G)\setminus H,
 \qquad \epsilon(v)=5-d(v).
\]
Thus $\epsilon(v)\geq0$ throughout the remaining argument.

\begin{lemma}\label{lem:local}
For every vertex $x$, let $r_x$ be the number of vertices $z\ne x$ with
$\codeg_G(x,z)=0$, and let $t_x$ be the number of triangles containing
$x$.  For $x\in H$, define
\[
 D_x=\sum_{y\in N(x)}\epsilon(y)
\]
Then
\begin{align}
 D_x&\geq4,\label{eq:D-lower}\\
 r_x&=D_x-4.\label{eq:r-exact}
\end{align}
For every vertex $x$, moreover,
\begin{equation}\label{eq:triangle-lower-local}
 r_x\geq d(x)-2t_x.
\end{equation}
In particular, $r_x\geq1$ whenever $d(x)$ is odd.
\end{lemma}

\begin{proof}
For fixed $x$, the sets
\[
 N(y)\setminus\{x\},\qquad y\in N(x),
\]
are pairwise disjoint.  Otherwise, two neighbors of $x$ and an element in
the intersection form a $4$-cycle with $x$.  They are subsets of
$V(G)\setminus\{x\}$, and hence
\begin{equation}\label{eq:local-packing}
 \sum_{y\in N(x)}(d(y)-1)\leq16.
\end{equation}
When $d(x)=5$, the left side equals $20-D_x$, proving
\eqref{eq:D-lower}.  Its disjoint union is exactly the set of vertices
$z\ne x$ that have a common neighbor with $x$.  Therefore
\[
 16-r_x=20-D_x,
\]
which is~\eqref{eq:r-exact}.

No edge can belong to two triangles, since the two third vertices would
form a $4$-cycle with the endpoints of that edge.  Consequently the
triangles through $x$ use $2t_x$ distinct edges incident with $x$.  The
other $d(x)-2t_x$ incident edges $xz$ lie in no triangle, so $x$ and each
such $z$ have no common neighbor.  This proves
\eqref{eq:triangle-lower-local}.  If $d(x)$ is odd, at least one incident
edge remains.
\end{proof}

We now sum these local quantities.  Define the low--low deficiency cost
\[
 C=\sum_{uv\in E(G[L])}\bigl(\epsilon(u)+\epsilon(v)\bigr).
\]
Counting low-to-high incidences, first from $H$ and then from $L$, gives the
exact identities
\begin{align}
 \sum_{x\in H}D_x&=S-C,\label{eq:sum-D}\\
 \sum_{x\in H}r_x&=S-C-4|H|.\label{eq:sum-r}
\end{align}

In cases 1--6 of Table~\ref{tab:degrees},
\eqref{eq:D-lower} and~\eqref{eq:sum-D} require
\[
 4|H|\leq S-C\leq S.
\]
But the respective pairs $(S,4|H|)$ are
\[
 (26,48),(26,44),(34,44),(34,44),(34,40),(34,36),
\]
a contradiction in every case.

In case 7, equation~\eqref{eq:sum-r} gives
\[
 \sum_{x\in H}r_x=42-C-36=6-C\leq6.
\]
All nine vertices of $H$ have odd degree, whereas
Lemma~\ref{lem:local} makes the same sum at least $9$.  This excludes the
seventh row.

\section{The final degree multiset}

It remains to exclude case 8.  Here $|H|=8$, while $L$ contains six
degree-$4$ vertices, two degree-$3$ vertices, and one degree-$2$ vertex.
Thus $S=42$, and~\eqref{eq:sum-r} becomes
\begin{equation}\label{eq:RH}
 R_H:=\sum_{x\in H}r_x=10-C.
\end{equation}
Every vertex of $H$ has odd degree, so $R_H\geq8$ and $C\leq2$.
Every low--low edge contributes at least $2$ to $C$, with equality only for
an edge between two degree-$4$ vertices.  It follows that
\begin{equation}\label{eq:q}
 q:=e(G[L])\in\{0,1\},\qquad C=2q.
\end{equation}

The degree sum in $L$ is $32$, so the number of high--low edges is
$32-2q$.  The degree sum in $H$ is $40$; consequently
\begin{equation}\label{eq:eH}
 e(G[H])=\frac{40-(32-2q)}2=4+q.
\end{equation}
Let
\[
 T_H=\sum_{x\in H}t_x
\]
be the number of high-vertex incidences with triangles.  Summing
\eqref{eq:triangle-lower-local} over $H$ and using
\eqref{eq:RH}--\eqref{eq:q} gives
\begin{equation}\label{eq:TH-lower}
 T_H\geq\frac{40-R_H}{2}=15+q.
\end{equation}

There is an incompatible upper bound.  Recall that the triangles of $G$
are edge-disjoint.  Charge their high-vertex incidences as follows.  A
triangle with two high vertices charges both high incidences to its unique
high--high edge.  A triangle with three high vertices charges its three
incidences to its three different high--high edges.  A triangle with one
high vertex charges that incidence to its unique low--low edge.  Thus each
high--high edge receives at most two charges and each low--low edge at most
one, whence, by~\eqref{eq:q} and~\eqref{eq:eH},
\begin{equation}\label{eq:TH-upper}
 T_H\leq2e(G[H])+e(G[L])=8+3q.
\end{equation}
For $q=0$ or $q=1$, the lower bound $15+q$ is larger than the upper bound
$8+3q$.  The contradiction eliminates case 8.  Lemma~\ref{lem:eight} is
exhaustive, so Theorem~\ref{thm:main} follows.

\section*{Computational and AI assistance disclosure}

Three exact Python enumerators were used for Lemma~\ref{lem:eight}.  Two
implement their graphicality tests directly; the third uses NetworkX~3.4.2
as an independent cross-check.  The complete programs and saved outputs
record the iteration ranges, predicates, and counts.  OpenAI Codex assisted
with literature searching, proof development, implementation and checking
of the exact enumerators, adversarial review organization, and manuscript
preparation.  Language-model output and successful program execution were
not treated as proofs by themselves: the mathematical argument is presented
in full, and the computational claim is restricted to the explicitly
exhausted finite degree-multiset domain.  Final responsibility for the
manuscript remains with the author.  No floating-point computation or proof
assistant was used.

\bibliographystyle{plain}
\bibliography{references}

\end{document}
